% Part: lambda-calculus
% Chapter: syntax
% Section: term-revisited

\documentclass[../../../include/open-logic-section]{subfiles}

\begin{document}

\olfileid{lam}{syn}{tr}
\olsection{Suku minangka Kelas Ekuivalensi-$\alpha$}

Sabanjure kita nganggep suku
nganti ekuivalensi-$\alpha$.
Tegese, nalika nulis sawijining
suku, kang dimaksud iku kelas
ekuivalensi-$\alpha$ kang ngemot suku
kasebut. Umpamane, kita nulis
$\lambd[a][\lambd[b][a c]]$
kanggo himpunan kabeh suku kang
ekuivalen-$\alpha$ karo suku iku,
kayata $\lambd[a][\lambd[b][a
    c]]$, $\lambd[b][\lambd[a][b c]]$,
lan sapiturute.

Ing bagean sadurunge, aksara
kayata $N, Q$ nglambangake suku;
sabanjure aksara iku nglambangake
kelas. Kelas-kelas kasebut, dudu
wakil suku tartamtu, kang bakal
dadi objek panliten sabanjure.
Aksara kayata $x, y$ tetep
nglambangake variabel.

Kita uga nggunakake notasi
$\rep{M}$ kanggo nglambangake
!!{element} sembarang saka kelas~$M$,
lan $\rep{M}[0]$, $\rep{M}[1]$,
lan sapiturute yen mbutuhake
luwih saka siji wakil.

Kanggo nyederhanakake andharan,
kita nggunakake maneh notasi
saka suku. Ing kelas ana
definisi iki:
\begin{defn}
  \begin{enumerate}
  \item $\lambd[x][N]$ ditetepake
    minangka kelas kang ngemot
    $\lambd[x][\rep{N}]$.
  \item $PQ$ ditetepake minangka
    kelas kang ngemot $\rep{P}\rep{Q}$.
  \end{enumerate}
\end{defn}

Definisi iki konsisten amarga
konversi-$\alpha$ kompatibel.

\begin{defn} \ollabel{def:fv}
  \emph{Variabel bebas} saka kelas
  ekuivalensi-$\alpha$~$M$, yaiku
  $\FV{M}$, ditetepake minangka
  $\FV{\rep{M}}$.
\end{defn}

Definisi iki konsisten amarga
$\FV{\rep{M}[0]} =
\FV{\rep{M}[1]}$, kaya kang
dituduhake ing \olref[alp]{thm:fv}.

Notasi substitusi uga digunakake
maneh kanggo kelas:
\begin{defn} \ollabel{def:sub}
  \emph{Substitusi} $R$ kanggo
  $y$ ing $M$, yaiku
  $\Subst{M}{R}{y}$, ditetepake
  minangka kelas saka
  $\Subst{\rep{M}}{\rep{R}}{y}$
  kanggo wakil $\rep{M}$
  lan $\rep{R}$ sembarang
  kang ndadekake substitusi
  kasebut katetepake.
\end{defn}

Definisi iki uga konsisten,
kaya kang dinyatakake dening
\olref[alp]{cor:sub}.

Gatekna yen definisi iki
nyederhanakake nalar kita.
Umpamane:
\begin{align}
  \Subst{\lambd[x][x]}{y}{x} & \ollabel{eq:1}\\
  &= \Subst{\lambd[z][z]}{y}{x} \ollabel{eq:2}\\
  &= \lambd[z][\Subst{z}{y}{x}] \\
  &= \lambd[z][z]
\end{align}

\olref{eq:1} wis katetepake
uga minangka substitusi suku
miturut aturan pangiket kang padha:
asile suku kasebut dhewe.
Ing kene substitusi dimangerteni
ing kelas, mula \olref{eq:2}
ngidini kita ngganti
$\lambd[x][x]$ nganggo wakil
$\lambd[z][z]$, amarga
loro-lorone ana ing kelas kang padha.

Kanthi alesan padha,
sabanjure kita nganggep wakil
kang dipilih tansah nyukupi
syarat substitusi. Umpamane,
nalika nemoni
$\Subst{\lambd[x][N]}{R}{y}$,
kita nganggep wakil
$\lambd[x][N]$ dipilih
saengga $x \neq y$
lan $x \notin \FV{R}$.

Amarga rada aneh nyebut
$\lambd[x][x]$ minangka
``kelas'', sabanjure kita
nyebut objek iki
suku-$\Lambda$
(utawa mung ``suku'' ing sisa
bagean iki), supaya beda
saka suku-$\lambda$ kang
wis kita kenal.

\begin{editorial}
  Kita durung bisa ninggalake
  suku lumrah: sakabehe definisi
  suku-$\Lambda$ adhedhasar
  suku-$\lambda$, lan kita durung
  menehi cara kanggo netepake
  fungsi ing suku-$\Lambda$.
  Mula fungsi kaya mangkono
  kudu luwih dhisik ditetepake
  ing suku-$\lambda$, banjur
  ``diproyeksikake'' menyang
  suku-$\Lambda$, kaya ing substitusi.
  Nanging kita nganggep pamaca
  bisa mangerteni kanthi intuisi
  carane fungsi ing suku-$\Lambda$
  bisa ditetepake.
\end{editorial}
\end{document}
